\subsection*{\emph{Spanned}：数据与缓冲区}

tileflow 程序的主要任务，就是搬动大规模数据集来处理它们。在某些术语体系里，Croqtile 函数的输入和输出分别称为「输入数据」和「输出数据」，其余存储则称为「缓冲区」。从函数的视角看，这就在概念上分出了「外部」内存和「内部」内存。不过「数据」和「缓冲区」指的都是存储位置；在 Croqtile 中，两者都定型为\textbf{带形状类型}，表明它们是与形状关联的数据或缓冲区。

这两种情况都称为带形状数据。\co@mdspan@ 描述形状，带形状对象则将形状、元素类型和存储访问组合起来。单独命名或推导一个形状，并不会复制数据。
